\chapter{What Counts as Knowledge?}
\section{A Stopped Watch}
Pick up the old watch at the back of a drawer.

Perhaps your grandfather left it. Its glass is scratched and its leather strap stiff. Crucially, it stopped long ago, its hands at 4:23, morning or afternoon; nobody remembers the day it stopped.

Try a peculiar thought experiment. Though dead, the watch is still right twice daily: at 4:23 in the morning and afternoon, its display matches the actual time exactly. Suppose you glance at it one afternoon and think it is 4:23, just as the real time happens to be 4:23.

Do you \textbf{know} it is 4:23?

Your belief is true, exactly. You did not guess blindly: you checked a watch, something you ordinarily trust. Yet almost everyone shakes their head. That is luck, not knowledge.

Strange. If a belief is true and apparently supported, why is it not knowledge? What separates knowing from happening to be right? That missing piece is our quarry. Rather than memorising definitions, we shall dismantle I know, inspect its components and ask why their assembly can still occasionally produce something unlike knowledge.

\section{On the Sand in Athens}
Enter Athens in the late fifth century BCE, at the home of a wealthy young man called Meno.

Afternoon sun slants across a courtyard smelling of olive oil and dust. An old man crouches with a stick, drawing a square in fine sand. He is Socrates, Athens's most famous and irritating questioner. Meno stands opposite, with an unnoticed young slave beside him, a child who has never attended school.

Meno and Socrates have been arguing over whether virtue can be taught. Meno poses a difficult objection: how can someone investigate what they do not know? Ignorance prevents knowing where to look; if they find it, how will they recognise it?

Socrates does not answer directly. He points to the child and proposes an experiment. A square with side two has area four: what side would double its area? The boy immediately says four, doubling the side. Socrates draws it, and the child sees the area is sixteen, quadrupled. Scratching his head, he suggests three. Drawing reveals nine, still wrong.

Notice what follows: Socrates \textbf{never gives} an answer. He draws and asks whether the child is sure, and what the calculation yields. Sweat appears on the boy's forehead; household slaves stop working to watch. Finally, studying the diagonal, he suddenly gives the answer: use the original square's diagonal as the new side, doubling the area.

Socrates tells Meno, in effect: he began ignorant and reached the right answer, not placed in his mouth by another person but drawn from his mind through successive questions.

Plato's Meno then supplies a more everyday illustration. Imagine two guides to Larissa. One has never visited but guesses the route correctly; the other has walked it. Both get you there, yet you distinguish knowing from guessing rightly. The first has true opinion; the second has knowledge.

Through Socrates, Plato says true opinion becomes knowledge when tethered by an account of its cause. We shall repeatedly return to test that formulation.

\section{What Is the Extra Piece?}
Set out the conflict without rushing to an answer.

Meno's difficulty is real. If true opinion and knowledge guide action equally well, getting us to Larissa or correct exam answers, why value knowledge more? Someone reproducing a memorised answer or guessing correctly scores as highly as someone understanding. Where is the difference between equally marked papers?

You might say real understanding. But what is understanding? Beware a circle: knowledge means real understanding, which means knowledge. To escape, dismantle knowing into components.

For over two millennia, philosophers' main recipe has required three simultaneous conditions:

\begin{itemize}
\item \textbf{You believe it.} You cannot know what you disbelieve. Reciting Earth's orbit around the Sun while inwardly believing the reverse does not amount to knowing the former.
\item \textbf{It is true.} Nobody can know a falsehood. You can firmly believe Earth is flat, but cannot know it, because it is not.
\item \textbf{You have adequate reasons.} Not a guess or coin toss, but evidence, inference or a dependable source.
\end{itemize}
Belief, truth and reasons: an elegant, clear three-part recipe scarcely altered for two thousand years. The lucky guide lacks the third: no reasons, only luck.

But recall the stopped watch. Its observer believes, satisfying one condition; says something true, satisfying another; and has a reason, consulting an ordinarily reliable watch, seemingly satisfying the third. All three are present, yet we hesitate.

Here is the real question: \textbf{does this ancient recipe omit something?} What would the missing component look like? More fundamentally, why assume our countless everyday claims to know deserve that description?

\section{Five Wells of Knowing}
Before repairing the recipe, step back to examine how information enters our heads. Two genuinely opposed positions deserve full hearings.

\textbf{First: I trust only what I see myself.}

You probably know such people, perhaps including yourself. Online news may be invented, books outdated, teachers mistaken. Only personal seeing and touching count.

Present their reasons fairly, rather than dismissing them as argumentative. Each person between an event and your ear adds an opportunity for error. Whispering games transform a sentence after ten people. Rumours, fake news and inherited misconceptions expose the chain's failures. Plenty of textbook facts have later been crossed out. These people set a high price on trust, not merely seek quarrels.

\textbf{Second: without others, you know almost nothing.}

Ask them: did you witness your own birth? Why know your birthday? Have you seen Qin Shi Huang, measured Everest or personally verified Earth's orbit?

No. Birthdays come from parents; Qin Shi Huang from texts and archaeology; Everest from surveyors; heliocentrism from centuries of astronomers. Inventory your knowledge and \textbf{over ninety per cent came through others}: parents, teachers, books, news and experts. Strictly restricting it to personal sight shrinks the store to a small room's touched objects and a few acquaintances. Consistently applied, the rule cannot even establish your country, since borders are not something you can see wherever you walk.

Each extreme has a fatal weakness. Eyewitness-only admission costs too much for anyone to enter; believing everything admits every fraud. Maturity lies in \textbf{keeping accounts of sources}: from which well did this claim come, and how has its water quality been recently?

Now inventory five wells:

\begin{itemize}
\item \textbf{Perception}: direct sight and hearing, quick but deceptive. Chopsticks look bent in water; trees seem to retreat from a speeding train.
\item \textbf{Memory}: stored perception retrieved later. Not a video but a retold story, it may blend childhood experiences with later photographs and others' accounts.
\item \textbf{Testimony}: everything others tell you, overwhelmingly humanity's main source, whose every link can break.
\item \textbf{Inference}: conclusions beyond direct observation. You have not seen the Sun's far side, but reason tells you Earth rotates. Detectives did not witness a crime, but fingerprints and surveillance can identify its perpetrator.
\item \textbf{Expert knowledge}: testimony's upgraded version, supplied by trained people checked by peers. More reliable, it still raises questions about who counts as expert and what to do when experts disagree.
\end{itemize}
This inventory is not modern. Around two millennia ago, India's Nyāya school, known for logic and debate, examined means of knowledge, pramāṇa: perception, inference, comparison, such as recognising a wild ox after being told it resembles a domestic cow, and reliable testimony. Other schools accepted two or proposed six, arguing for centuries. They debated not superiority among kinds of knowledge, but \textbf{how many legitimate entrances exist}: admit no extra gate for impostors, close none needed by genuine knowledge.

Beware another mistaken ranking: perception first, testimony last. Forensic psychology offers a reversal. Confident eyewitness identification is notoriously error-prone; memory repairs scenes into plausible shapes, and confidence often fails to track accuracy. Cross-checked indirect information, such as weather-radar data, can outperform eyes looking through rain. \textbf{Sources' rankings change with circumstances}. What matters is not the well's name, but its likelihood of error here.

\section[Thinking Bridge: Three Gates]{Thinking Bridge: Three Gates for ``I Know''}
Assemble a portable tool. Before saying I know, or evaluating someone else's claim that you should know, send it through three gates:

\textbf{First: do you believe it?}

Knowledge must be your own position, not a visitor passing through your mouth. Full exam marks alongside complete disbelief leave a student empty in knowledge's strict sense. This gate excludes repetition.

\textbf{Second: is it true?}

Hardest and most troublesome, this gate is controlled by the world, not you. Firm belief and excellent reasons cannot make knowledge if reality differs. Displaced theories such as phlogiston and geocentrism once combined conviction and strong reasons, but fail this test. It excludes sincere error.

\textbf{Third: do your reasons qualify?}

Lucky guesses are not knowledge. Reasons need a genuine route to conclusions through evidence, inference or reliable sources. This gate excludes luck.

Add a neglected scale: \textbf{reasons have assessable strength}. Vaguely remembering red clothing and surveillance showing it pass the same gate at different levels. There is an intermediate state, \textbf{reasonable belief}: insufficient grounds to assert knowledge confidently, but enough for informed action. A seventy-per-cent rain forecast does not establish tomorrow's rain as known, yet justifies an umbrella. Knowledge demands full strength; reasonable belief permits a discount. Ninety per cent of everyday supposed knowledge lives in that discounted area. Acknowledging it begins honesty, not shame.

Return to Athens. The lucky guide passes belief and truth, but has no supported route through the third gate. The experienced guide passes all three. The recipe looks complete until the stopped watch appears.

\section{Two Ancient Statements of Honesty}
Read primary sources. Chinese thinkers left two brief, forceful statements about knowing over two millennia ago.

In the Analects' Governing chapter, Confucius tells Zilu, whose personal name was Zhong You:

\begin{quote}
The Master said, ``You, shall I teach you what knowing is? Recognise what you know as known, and what you do not know as unknown. That is knowing.''
\end{quote}
In brief: Zhong You, call knowledge knowledge and ignorance ignorance; that is genuine knowing.

Notice the oddity: knowing your ignorance counts as knowledge. This is no word game. Confucius adds a reverse inspection: not only whether beliefs passed the gates, but whether failed items entered the successful queue. An honestly marked gap makes someone epistemically richer than mixing guesses with knowledge, because their inventory is accurate. Later scholars call this second-order knowledge, knowledge about one's knowledge. Socrates similarly questioned dignitaries claiming expertise in courage or justice until uncertainty emerged, then located his advantage in knowing his ignorance. Distant civilisations, without correspondence, reached a similar position around the same era: inventory ignorance to begin knowledge.

The second source is Autumn Floods in Zhuangzi, the debate with Huizi above the Hao:

\begin{quote}
Zhuangzi and Huizi strolled above the Hao. Zhuangzi said, ``The minnows swim freely at ease. That is fishes' happiness.'' Huizi said, ``You are not a fish. How do you know their happiness?'' Zhuangzi replied, ``You are not me. How do you know I do not know their happiness?''
\end{quote}
They walk on the bridge; Zhuangzi calls the leisurely minnows happy. Huizi asks how a non-fish knows. Zhuangzi asks how someone not himself knows his ignorance.

Often treated as amusing quibbling, Huizi's question targets our third gate: \textbf{where do your reasons come from?} Zhuangzi lacks fishes' sensory channels; fish offer no verbal testimony. Which well supplies happiness? Asking whether AI feels has the same structure. Zhuangzi's response is not mere evasion: Huizi's standard rebounds. If not being another prevents knowing, Huizi cannot know whether Zhuangzi knows, undermining his challenge. In a few exchanges, they explore the still-unsettled problem of other minds.

Together, Confucius addresses \textbf{inward} honesty about knowledge and ignorance; Zhuangzi and Huizi address \textbf{outward} access to another mind. New vocabulary over millennia has bypassed neither difficulty.

\section{The Stopped-Watch Puzzle: Three Accounts}
Return directly to the watch stopped at 4:23.

In 1963, Edmund Gettier published a three-page paper. Since it is not public-domain text, we paraphrase its central idea. He constructed situations meeting belief, truth and adequate reasons while almost everyone judged them not knowledge. The stopped watch typifies such cases: all gates appear passed, but the observer hits a twice-daily coincidence. These became Gettier problems, a needle bursting an ancient balloon. The three conditions are necessary, not sufficient; passing all need not produce knowledge.

Why? Several genuinely different explanations address these facts.

\textbf{First: something false entered the reasons.}

The chain is: the watch displays 4:23$\rightarrow$the watch works normally$\rightarrow$therefore it is 4:23. Its middle link is \textbf{false}. This repair requires no false premises in knowledge-producing reasoning. It diagnoses the watch precisely, but our everyday chains are long and hidden. Who can verify every link? Knowing water boils at one hundred Celsius depends on teachers, books, experiments and accurate instruments. Check rigorously and few claims remain safe to sign.

\textbf{Second: the problem concerns origins, not reasons.}

Examine \textbf{how belief is produced}. A working watch is a reliable production line for time beliefs; a stopped one is defunct, despite an occasional good product. Reliabilism calls knowledge true belief produced reliably. This liberates knowers from articulating every reason and accommodates children's and animals' reliable perception. Yet consider a county full of cardboard barn façades and one real barn. Driving past, you happen to see the real one with perfectly functioning eyes. Turning a second earlier would have shown a façade. Do you know there is a barn? Most still hesitate. Reliabilism must specify reliability and its relevant scope, reopening debate.

\textbf{Third: perhaps knowledge has no recipe.}

Endlessly layered repairs may indicate a wrong approach. Knowing the route to a friend's home, knowing their unhappiness and knowing Pythagoras involve different ordinary uses. One three-part recipe may be no more plausible than one key for every lock. Gettier cases disturb us because our intuitions mix different standards. A bracing view, but costly: courts must judge informed witnesses; science must distinguish known and unknown. How without a common standard? Abandoning a recipe is easier than managing the consequences.

Each account grasps part, not all, of the truth. Here is an important result: \textbf{knowledge, a word used every minute, still lacks a universally accepted complete definition}. This reflects not philosophical incompetence but depth, at the intersection of belief, world and reasons. A shift in any requires rebuilding.

\section[Further Challenge: If Everything Is a Dream]{Further Challenge: What If Everything Is a Dream?}
Push doubt to its limit and inspect knowledge's foundations with Descartes.

In Meditations on First Philosophy, he describes sitting by a winter fire in a dressing gown, paper in hand. Why be certain? Dreams have presented the same warmth and scene, later revealed unreal. Can he \textbf{prove} he is not dreaming now? He finds no way. Dream experience can perfectly mimic waking; wakefulness cannot verify itself from within.

Go deeper: might a powerful demon systematically supply false beliefs? Senses deceive; even arithmetic may be unsafe. Perhaps the demon makes two plus three equalling five merely feel obvious. In our terms, the truth gate has no directly accessible key. Experience is glass between us and the world, whose external existence and agreement with its image cannot be inspected from inside.

Others took this path. About five centuries earlier, Persian scholar al-Ghazālī's Deliverance from Error describes a crisis of distrust in perception and reason lasting months before recovery. Earlier still comes Zhuangzi's butterfly dream. Zhuang Zhou dreams himself a happily fluttering butterfly, unaware of himself, then wakes as Zhou and wonders whether Zhou dreamed butterflyhood or a butterfly now dreams Zhou. Playfulness and Descartes's fireside fear touch the same stone: \textbf{wakefulness lacks a certificate delivered from outside the dream}. Radical doubt recurs in human minds, not one culture alone.

Descartes's unbreakable stone is famous: Cogito, ergo sum, I think, therefore I am. Doubting everything cannot erase doubting itself, performed in the attempt. At least an I exists. From there he tried rebuilding knowledge brick by brick. Its stability has been disputed for three centuries, but the method exceeds the conclusion: \textbf{to understand knowledge, test how severe a doubt it survives}.

After Descartes, European philosophers divided over which well mattered most.

\textbf{Empiricism}, represented by Locke and Hume, emphasised experience. Locke's blank-slate metaphor denies innate ideas, attributing mental contents to experience. Hume dismantled causation: ten thousand sunrises encourage belief in tomorrow's, but no logical guarantee connects past recurrence to future recurrence. Habit, not proof, sustains expectation. Every inferential step from observed to unobserved rests on experiential soil rather than logical steel.

\textbf{Rationalism}, including Descartes, drew confidence from mathematics. A triangle's angles sum to 180 degrees not because a thousand measurements say so; measurements contain error while the theorem is absolute, eternal and universal. Such knowledge seems independent of experience, even commanding it.

After over a century, Kant supplied a synthesis. Hume, he acknowledged in the Prolegomena to Any Future Metaphysics That Will Be Able to Come Forward as Science, awakened him from dogmatic slumber. In modern terms, empiricists make knowledge \textbf{imported goods}, rationalists \textbf{home production}; Kant requires both. The world supplies sensory materials, while the mind supplies processing machinery: time, space and causation as factory settings. Materials must be processed; machinery without them runs empty. Tomorrow's sunrise is neither mere guess nor pure proof, but their combined product, reliable enough for science, certified by mental machinery rather than the universe's own signature. Debate continues over this synthesis, but its introductory lesson is clear: \textbf{asking where knowledge comes from eventually asks what the knower is like}.

\section{When Pictures No Longer Prove It}
These ancient questions recur in your pocket at greater difficulty.

\textbf{Failure at gate one}: your interlocutor may not be human. An AI model might answer a hundred questions correctly. Does it believe them? Inside it is calculation and output, not a position of thinking something so. Our recipe stops it at belief: truth is frequent and training may count as a kind of support, but belief is absent. This extends Gettier into life: can something always right but never believing know? Your answer shapes trust in machines.

\textbf{Failure at gate two}: visual testimony depreciates. Pictures once ended online arguments; face-swapped videos and generated photographs now expose direct sight to systematic poisoning for the first time. India's long debate over perception's reliability becomes our technical question: was this video filmed or algorithmically drawn?

\textbf{Failure at gate three}: conspiracy theories imitate chains of reasons. A carefully fabricated theory can be internally coherent and apparently supported throughout. Its illness lies in the first repair's diagnosis: a deeply buried false premise, armoured with accusations that doubt proves naivety. Ask not whether it hangs together, but whether each link withstands independent checking.

Degrees matter too. Phones show seventy-per-cent precipitation, ninety-five-per-cent clinical effectiveness and ninety-nine-per-cent detection confidence. Science has abandoned postures of complete certainty for scales. Mature knowledge permits graded belief rather than total acceptance or rejection: \textbf{ask how large the discount is, then what justifies it}. Explaining seventy-per-cent confidence comes closer to knowledge than constant certainty. Confucius anticipated this long ago.

\section{Your Turn: Who Sets the Standard?}
Return to the watch and lay out our account.

A stopped watch happens to show the correct time, producing belief, truth and reasons. Does its observer know? Your intuition probably says no. Now answer the final question with reasons:

\textbf{Most everyday claims to know cannot survive scrutiny at this level}. Tomorrow's mathematics lesson rests on a timetable that might change. A friend's address rests on testimony that might misremember a number. Under strict philosophical standards, you, I and everyone possess little daily knowledge, mostly reasonable beliefs that happen not to go wrong.

What should we do? Reserve know for near-certainty and admit mainly reasonable belief? Allow everyday knowing some luck, leaving laboratories their exacting standards? Or choose a third route: standards depend on context rather than philosophers alone, varying between classroom, courtroom, operating theatre and conversation?

There is no standard answer. Remember the recurring stone: every I know is a signature warranting passage through three gates. Before speaking, pause half a second and count yours. That half-second is this chapter's intended gift.
